\documentclass[a4paper,12pt, french, noSujetB]{article}

\usepackage{layout_evaluation}

\classe{Seconde}
\titre{Concours des médianes}
\evalnumber{2}
\evaltype{Exercice autonome}
\date{septembre 2026}

\setlist{leftmargin=5mm}

\begin{document}
	On veut prouver que les médianes d'un triangle sont concourantes. Plus précisément, elles se coupent à $\frac{2}{3}$ de la longueur de la médiane.

	On se donne donc un triangle ABC quelconque. On note I la base de la médiane partant de C, J la base de la médiane partant de A et K la base de la médiane partant de B.

	\begin{enumerate}
		\item Faire un dessin.
		\item Compléter l'égalité suivante pour traduire la phrase \og elles se coupent à $\frac{2}{3}$ de la longueur de la médiane\fg{} par une égalité de vecteurs:
		$$\frac{2}{3}\vec{AJ} = \vec{AB}+\frac{2}{3}\vec{BK} = \dots$$
		\item Montrons la première égalité. On veut décomposer tous les vecteurs en fonction de deux vecteurs choisis au préalable, qu'on appellera une \textbf{base}. Choisissons par exemple $\vec{AB}$ et $\vec{BC}$.
		\begin{enumerate}
			\item
			Exprimer $\vec{AJ}$ dans cette base.
			\item Exprimer $\vec{BK}$ dans cette base.
			\item Prouver l'égalité de la question 2.
		\end{enumerate}
	\end{enumerate}
\end{document}
